\documentclass[11pt]{article}
\usepackage[T1]{fontenc}
\usepackage{amsmath,amssymb,amsthm,lmodern}
\usepackage[margin=1in]{geometry}
\usepackage{hyperref}
\hypersetup{hidelinks,pdftitle={EGNO Proposition 4.2.10: from left duals to rigidity},pdfauthor={Pavel Etingof; Shlomo Gelaki; Dmitri Nikshych; Victor Ostrik}}
\DeclareMathOperator{\Hom}{Hom}
\DeclareMathOperator{\End}{End}
\DeclareMathOperator{\FPdim}{FPdim}
\DeclareMathOperator{\Gr}{Gr}
\newcommand{\Cat}{\mathcal C}
\newcommand{\unit}{\mathbf 1}
\newcommand{\Obj}{\mathcal O}
\newtheorem*{primaryprop}{Proposition 4.2.10}
\title{Tensor Categories\\Proposition 4.2.10: From Left Duals to Rigidity}
\author{Pavel Etingof \and Shlomo Gelaki \and Dmitri Nikshych \and Victor Ostrik}
\date{Private reader of the 2015 AMS author final copy}
\begin{document}
\maketitle

% BEGIN CONTEXT EGNO:4.2.10
\noindent\textbf{Fixed source and scope.}
\emph{Tensor Categories}, Mathematical Surveys and Monographs 205,
American Mathematical Society, 2015, author final copy.
\par\noindent\url{https://math.mit.edu/~etingof/egnobookfinal.pdf}.
\par\noindent
The retained PDF has 362 physical pages; its metadata record creation on
May 13, 2015 and modification on July 29, 2015. The complete primary is
on physical pages 84--85, printed pages 67--68. Lemma 4.2.11 and the
paragraph after its proof belong to this one primary proof.

The chapter opens on physical page 82 (printed 65) with ``Let $k$ be an
algebraically closed field.'' The hypothesis here is a \emph{finite ring
category with left duals}; right duals are the conclusion. No
characteristic-zero, semisimplicity, symmetry, or identification of the
double left dual with the identity is added. Selected earlier and forward
same-book inputs are printed below as uncounted support. In particular,
Proposition 3.3.13(1) is an explicit strong input. Proposition 4.5.7,
whose hypothesis already says tensor categories, is not used in this
reader to obtain the missing right duals.

\noindent\textbf{Source notice.}
The source footer states: ``Author's final version made available with
permission of the publisher, American Mathematical Society. License or
copyright restrictions may apply to redistribution; see
\url{http://www.ams.org/publications/ebooks/terms}.'' This reader remains
private. The separate February 14, 2026 corrections are identified at the
end; they do not silently replace the 2015 text.

\subsection*{Terminology and explicit prior inputs (uncounted)}
\paragraph{Finite abelian categories.}
Definition 1.8.5 (physical 26, printed 9) says that a $k$-linear abelian
category $\Cat$ is \emph{finite} if it is equivalent to the category
$A\text{-mod}$ of finite dimensional modules over a finite dimensional
$k$-algebra $A$. Definition 1.8.6 gives the equivalent intrinsic conditions:
\begin{enumerate}
\item $\Cat$ has finite dimensional spaces of morphisms;
\item every object of $\Cat$ has finite length;
\item $\Cat$ has enough projectives, i.e., every simple object of $\Cat$
has a projective cover; and
\item there are finitely many isomorphism classes of simple objects.
\end{enumerate}
The first two conditions define \emph{locally finite} (Definition 1.8.1).
The notation $\Obj(\Cat)$ denotes the set of isomorphism classes of
simple objects. Length means Jordan--H\"older length. These definitions
do not identify the monoidal product with the standard tensor product
of modules or supply a fiber functor.

Definition 1.6.5 (physical 23, printed 6) calls $P$ projective if
$\Hom_{\Cat}(P,-)$ is exact. Under the standing finite-length condition,
Definition 1.6.6 calls an epimorphism $p:P(X)\to X$ a \emph{projective
cover} if $P(X)$ is projective and, whenever $g:P\to X$ is an
epimorphism from a projective object, there is an epimorphism
$h:P\to P(X)$ with $ph=g$. If a projective cover exists, it is unique
up to a non-unique isomorphism. The discussion on physical 27 (printed
10) constructs a projective generator from the covers of the finitely
many simple objects. That discussion is background, not a new target.

\paragraph{Ring and tensor categories.}
Definition 4.2.3 (physical 83, printed 66) reads:
A \emph{multiring category} over $k$ is a locally finite $k$-linear
abelian monoidal category $\Cat$ with bilinear and biexact tensor product.
If in addition $\End_{\Cat}(\unit)=k$, we will call $\Cat$ a
\emph{ring category}.

Definition 4.1.1 (physical 82, printed 65) starts with a locally finite
$k$-linear abelian \emph{rigid} monoidal category and calls it a
multitensor category when the tensor product is bilinear on morphisms;
with $\End_{\Cat}(\unit)\cong k$ it is a tensor category.
Definition 2.10.11 (physical 59, printed 42) calls an object rigid if it
has left and right duals, and a category rigid if every object is rigid.

\paragraph{Orientations of duals (Definitions 2.10.1--2.10.2).}
On physical 57 (printed 40), a left dual $X^*$ has evaluation and
coevaluation
\[
\operatorname{ev}_X:X^*\otimes X\longrightarrow\unit,
\qquad
\operatorname{coev}_X:\unit\longrightarrow X\otimes X^*,
\]
and a right dual ${}^*X$ has
\[
\operatorname{ev}'_X:X\otimes{}^*X\longrightarrow\unit,
\qquad
\operatorname{coev}'_X:\unit\longrightarrow{}^*X\otimes X.
\]
Their two triangle identities are the definition's (2.43)--(2.46).
The preceding coherence convention (Remark 2.9.3) permits suppressing
brackets and the associated canonical associativity and unit maps when
there is no confusion. These definition excerpts specify the orientations;
the four triangle composites and the dual-morphism formulas (2.47)--(2.48)
are background and are not reprinted here.

Exercise 2.10.7(b) (physical 59, printed 42) states: if $U,V$ have left
(respectively, right) duals, then $U\otimes V$ has left dual
$V^*\otimes U^*$ (respectively, right dual ${}^*V\otimes{}^*U$).
The exercise's verification is a stated prior input; no new solution is
inserted here.

\paragraph{Proposition 2.10.8(a), with its printed proof.}
Let $\Cat$ be a monoidal category and let $V$ be an object in $\Cat$.
If $V$ has a left dual $V^*$ then there are natural adjunction
isomorphisms
\begin{align}
\Hom_{\Cat}(U\otimes V,W)&\xrightarrow{\sim}
\Hom_{\Cat}(U,W\otimes V^*),\tag{2.49}\\
\Hom_{\Cat}(V^*\otimes U,W)&\xrightarrow{\sim}
\Hom_{\Cat}(U,V\otimes W).\tag{2.50}
\end{align}
\begin{proof}[Printed proof (physical 59, printed 42)]
An isomorphism in (2.49) is given by
$f\mapsto(f\otimes\operatorname{id}_{V^*})\circ
(\operatorname{id}_U\otimes\operatorname{coev}_V)$ and has the inverse
$g\mapsto(\operatorname{id}_W\otimes\operatorname{ev}_V)\circ
(g\otimes\operatorname{id}_V)$. Other isomorphisms are similar and
are left to the reader as an exercise.
\end{proof}
Only the left-dual part (a) of the statement is selected. The printed
proof paragraph is retained in full, including its instruction concerning
the other isomorphisms; the right-dual part (b) is not assumed.
% END CONTEXT EGNO:4.2.10

% BEGIN SUPPORT EGNO:4.2.10 EXACTNESS
\subsection*{Exactness support: Proposition 4.2.9}
\noindent\emph{Physical 84, printed 67; complete statement and proof.}
If $\Cat$ is a multiring category with left duals, then the left
dualization functor is exact. The same applies to right duals.
\begin{proof}
Let $0\to X\to Y\to Z\to0$ be an exact sequence. We need to show
that the sequence $0\to Z^*\to Y^*\to X^*\to0$ is exact. Let $T$
be any object of $\Cat$, and consider the sequence
\[
0\to\Hom_{\Cat}(T,Z^*)\to\Hom_{\Cat}(T,Y^*)\to\Hom_{\Cat}(T,X^*).
\]
By Proposition 2.10.8, it can be written as
\[
0\to\Hom_{\Cat}(T\otimes Z,\unit)
\to\Hom_{\Cat}(T\otimes Y,\unit)
\to\Hom_{\Cat}(T\otimes X,\unit).
\]
Using the definition of a multiring category, we see that the last
sequence is exact. This implies that the sequence
$0\to Z^*\to Y^*\to X^*$ is exact.

Similarly, consider the sequence
\[
0\to\Hom_{\Cat}(X^*,T)\to\Hom_{\Cat}(Y^*,T)\to\Hom_{\Cat}(Z^*,T).
\]
By Proposition 2.10.8, it can be written as
\[
0\to\Hom_{\Cat}(\unit,X\otimes T)
\to\Hom_{\Cat}(\unit,Y\otimes T)
\to\Hom_{\Cat}(\unit,Z\otimes T),
\]
which is exact since $0\to X\otimes T\to Y\otimes T\to Z\otimes T$
is exact by the definition of a multiring category. This implies that
the sequence $Z^*\to Y^*\to X^*\to0$ is exact.
\end{proof}
% END SUPPORT EGNO:4.2.10 EXACTNESS

% BEGIN PRIMARY STATEMENT EGNO:4.2.10
\begin{primaryprop}
A finite ring category $\Cat$ with left duals is a tensor category
(i.e., it also has right duals).
\end{primaryprop}
% END PRIMARY STATEMENT EGNO:4.2.10

% BEGIN PRIMARY PROOF EGNO:4.2.10
\begin{proof}
We know that the left duality functor $X\mapsto X^*$ is fully faithful
(see Proposition 2.10.8), exact (see Proposition 4.2.9), and monoidal
(see Exercise 2.10.7(b)). In particular, this functor respects the
Frobenius--Perron dimension by Proposition 3.3.13(1). We claim that this
functor sends simple objects to simple ones. Let
$d_1<d_2<\cdots<d_r$ be the list of possible Frobenius--Perron dimensions
of simple objects in $\Cat$. Let $\Obj_i(\Cat)\subset\Obj(\Cat)$ be the
isomorphism classes of simple objects $L$ with $\FPdim(L)=d_i$,
$i=1,\ldots,r$.

\medskip\noindent\textbf{Lemma 4.2.11.}
\emph{Let $L\in\Obj_i(\Cat)$. Then $L^*$ is simple and there exists
$L'\in\Obj_i(\Cat)$ such that $L\simeq(L')^*$.}

\smallskip\noindent\emph{Proof.}
We prove the lemma by induction in $i$. The base ($i=1$) is obvious.
If $L^*$ is not simple then there exists a nonzero map $K\to L^*$,
where $K$ is simple with $\FPdim(K)<\FPdim(L^*)=d_i$. By the
induction assumption $K\simeq(K')^*$, where $K'$ is simple with
$\FPdim(K')<d_i$. Since the duality functor is full, we get a nonzero
map $L\to K'$, which is a contradiction. Thus we have a map
$L\mapsto L^*$ from $\Obj_i(\Cat)$ to itself. Again by the fullness of
duality we see that this map is injective. Since the set
$\Obj_i(\Cat)$ is finite, it is also surjective, which proves the lemma.
\hfill$\square$

\medskip
By Lemma 4.2.11 the functor $X\mapsto X^{**}$ induces a permutation
of $\Obj(\Cat)$. Thus there exists a finite power $F=(**)^{n}$ of the
double duality functor such that $F(L)\simeq L$ for any
$L\in\Obj(\Cat)$. Also, the functor $F$ is exact and fully faithful.
Let $P$ be the projective cover of $L\in\Obj(\Cat)$. Since the functor
$F$ is full we see that the object $F(P)$ has a unique simple quotient
$F(L)\simeq L$. It follows that we have a surjection $P\to F(P)$.
Since the length of $F(P)$ equals the length of $P$, this must be an
isomorphism $P\simeq F(P)$. Thus $F$ is surjective on isomorphism
classes of projective objects. Since any object of $\Cat$ is isomorphic
to a cokernel of a map of projective objects and the functor $F$ is
fully faithful, we get that $F$ is surjective on isomorphism classes of
objects and hence is an equivalence. Thus the duality functor
$X\mapsto X^*$ is surjective on isomorphism classes of objects. The
proposition is proved.
\end{proof}
% END PRIMARY PROOF EGNO:4.2.10

% BEGIN SUPPORT EGNO:4.2.10 GROTHENDIECK-FPDIM
\subsection*{Grothendieck ring and FPdim inputs (uncounted)}
\paragraph{Multiplication and Lemma 4.5.1.}
\emph{Physical 88--89, printed 71--72.}
In Section 1.5 the book introduced the Grothendieck group $\Gr(\Cat)$
of a locally finite abelian category. Now let $\Cat$ be a multiring
category over $k$. Let $X_i$, $i\in I$, be representatives of the
isomorphism classes of simple objects in $\Cat$. The tensor product
on $\Cat$ induces a natural multiplication on $\Gr(\Cat)$ defined by
\begin{equation}
X_iX_j:=[X_i\otimes X_j]
=\sum_{k\in I}[X_i\otimes X_j:X_k]X_k,\qquad i,j\in I.
\tag{4.7}
\end{equation}
Here $[X:Y]$ is the multiplicity of the simple object $Y$ in the
Jordan--H\"older series of $X$.

\noindent\textbf{Lemma 4.5.1.}
\emph{The above multiplication on $\Gr(\Cat)$ is associative.}
\begin{proof}
Since the tensor product functor is exact,
\[
[(X_i\otimes X_j)\otimes X_p:X_l]
=\sum_k[X_i\otimes X_j:X_k][X_k\otimes X_p:X_l],\quad\forall i,j,p,l.
\]
On the other hand,
\[
[X_i\otimes(X_j\otimes X_p):X_l]
=\sum_k[X_j\otimes X_p:X_k][X_i\otimes X_k:X_l].
\]
Therefore, the associativity of the multiplication follows from the
isomorphism $(X_i\otimes X_j)\otimes X_p\cong
X_i\otimes(X_j\otimes X_p)$.
\end{proof}
Thus, $\Gr(\Cat)$ is a $\mathbb Z_+$-ring with the unit $[\unit]$.
Definition 4.5.2 calls it the Grothendieck ring of $\Cat$.

\paragraph{Theorem 4.3.8(i): explicit statement input.}
\emph{Physical 86, printed 69.}
In a ring category with left duals, the unit object $\unit$ is simple.
Its complete printed proof is on physical 86--87 (printed 69--70);
that proof is available in the fixed source but is not reprinted here.
It is an explicit strong support statement for the following printed
proof, not a missing source or an assumption of right duals.

\paragraph{Proposition 4.5.4, complete statement and proof.}
If $\Cat$ is a ring category with left duals then $\Gr(\Cat)$ is a
transitive unital $\mathbb Z_+$-ring in the sense of Definition 3.3.1.
\begin{proof}
Recall from Theorem 4.3.8 that the unit object $\unit$ in $\Cat$ is
simple. So $\Gr(\Cat)$ is unital. This implies that for any simple
objects $X,Z$ of $\Cat$, the object $X\otimes X^*\otimes Z$ contains
$Z$ as a composition factor (as $X\otimes X^*$ contains $\unit$ as a
composition factor), so one can find a simple object $Y_1$ occurring in
$X^*\otimes Z$ such that $Z$ occurs in $X\otimes Y_1$. Similarly,
the object $Z\otimes X^*\otimes X$ contains $Z$ as a composition
factor, so one can find a simple object $Y_2$ occurring in $Z\otimes X^*$
such that $Z$ occurs in $Y_2\otimes X$. Thus $\Gr(\Cat)$ is transitive.
\end{proof}
In particular, Proposition 4.5.4 shows that for a ring category $\Cat$
with left duals, we can talk about Frobenius--Perron dimensions of
objects and regular objects, see Section 3.3.

\paragraph{Definitions 3.3.1 and 3.3.3 and prior character input.}
\emph{Physical 70, printed 53.}
For a $\mathbb Z_+$-ring $A$ with its distinguished basis $I$,
Definition 3.3.1 calls $A$ transitive if for any $X,Z\in I$ there exist
$Y_1,Y_2\in I$ such that $XY_1$ and $Y_2X$ contain $Z$ with a
nonzero coefficient. For a transitive unital $\mathbb Z_+$-ring of
finite rank, $\FPdim(X)$, $X\in I$, is the maximal non-negative
eigenvalue of left multiplication by $X$, and $\FPdim$ is extended
from $I$ by additivity. Definition 3.3.3 calls this function the
Frobenius--Perron dimension.

Proposition 3.3.6(1),(3) states that $\FPdim:A\to\mathbb C$ is a ring
homomorphism and is the unique character of $A$ taking non-negative
values on $I$, and that these values are actually strictly positive.
These are explicit strong prior inputs; their proof on physical 70--71
(printed 53--54), and the preceding Frobenius--Perron theorem, are not
reprinted here.

\paragraph{Proposition 3.3.13(1), with its printed proof.}
\emph{Physical 72, printed 55.}
Let $A_1,A_2$ be transitive unital $\mathbb Z_+$-rings of finite rank
with $\mathbb Z_+$-bases $I_1,I_2$. Let $f:A_1\to A_2$ be a unital
homomorphism, whose matrix in the bases $I_1,I_2$ has non-negative
entries. Then $f$ preserves Frobenius--Perron dimensions.
\begin{proof}[Printed proof of part (1)]
The function $X\mapsto\FPdim(f(X))$ is a character of $A_1$ with
non-negative values on the basis. By Proposition 3.3.6(3),
$\FPdim(f(X))=\FPdim(X)$ for all $X$ in $I$.
\end{proof}
Only part (1) and its own proof are selected. The source's final $I$
is retained; the statement specifies $I_1,I_2$. Parts (2),(3), concerning
regular elements, are not needed or reprinted.
% END SUPPORT EGNO:4.2.10 GROTHENDIECK-FPDIM

\subsection*{Separate 2026 corrections and attribution}
The six-page author correction document is dated February 14, 2026:
\url{https://math.mit.edu/~etingof/tensbookcor.pdf}.
Its physical page 3 says: ``Paragraph before Example 4.5.5: the ring
category should be assumed finite.'' The 2015 paragraph immediately
after Proposition 4.5.4 is retained above; this correction is recorded
separately. The present primary already assumes finiteness.
The same correction page changes the reference in Proposition 4.2.8
and interchanges the words ``monomorphisms'' and ``epimorphisms'' in
Corollary 4.3.9. Neither of those passages is used as a printed support
statement in this reader; they are not silently imported into the
primary or Proposition 4.2.9.

The exposition and proof collected here are by the four book authors.
The bibliographical notes on physical 105 (printed 88), Section 4.2,
specifically credit Proposition 4.2.1 to Deligne--Milne and Proposition
4.2.12 to Kazhdan--Lusztig and related sources. They do not give a
separate route byline for Proposition 4.2.10. Those neighbouring credits
are not reassigned to this proof. The original base-case, finite-stratum,
projective-cover, length and cokernel compressions remain at the authors'
printed level; no new proof is added.
\end{document}
